\section{Преглед на съществуващите решения}
\label{sec:existing_solutions}

Анализът на съществуващите решения е необходим, за да се оцени дали предложената архитектура действително допринася с нова стойност, или само преформулира вече познати практики. При C++ библиотеките за достъп до данни могат да се разграничат две независими оси: нивото на абстракция и моментът, в който се изгражда и валидира заявката. Нивото на абстракция варира от почти директно използване на клиентски C API до пълно ORM поведение. Моментът на валидиране може да бъде runtime, hybrid или compile-time.

Тези две оси са особено важни в контекста на настоящата дипломна работа. Ако се наблюдава само ``удобството на API-то'', лесно може да се пропусне съществената разлика между библиотека, която просто прави SQL низовете по-четими, и библиотека, която действително прехвърля структурната коректност към компилатора. Ако, от друга страна, се гледа само кога възниква грешката, може да се подцени цената на различните нива на абстракция, build toolchain сложност и пригодност за работа с повече от един storage модел.

\subsection{Класификация на подходите}

В най-ниския слой стоят native клиентските библиотеки като \code{SQLite C API}~\cite{sqlite_api}, \code{libpq}~\cite{libpq} и \code{MySQL Connector/C}~\cite{mysql_connector}. Те предоставят максимален контрол върху заявките, параметрите и резултатите, но не предлагат логическа абстракция. Приложението е принудено да работи директно със SQL низове, индекси на колони, позиции на параметри и backend-специфични типове.

Следващата група са runtime query builder библиотеките. Те обикновено подобряват четимостта на заявките и намаляват част от шаблонния код, но продължават да разчитат на runtime изграждане на текстово представяне. Типовата безопасност е частична и обикновено спира до границата на API повикването. Структурната валидност на заявката остава проблем на изпълнението.

Най-високото ниво на абстракция е ORM подходът. Там C++ класовете се асоциират с таблици, колони и връзки, а библиотеката генерира необходимите заявки. Предимството е намаляване на boilerplate кода. Недостатъкът е, че често се въвеждат външни генератори, макроси, анотации и opaque runtime метаданни, които затрудняват проследимостта между C++ кода и реалната заявка.

\begin{figure}[H]
\centering
\begin{tikzpicture}[x=2.0cm,y=1cm]
    \draw[->, thick] (0,0) -- (9,0) node[right] {\small момент на валидация};
    \draw[->, thick] (0,0) -- (0,7) node[above] {\small ниво на абстракция};

    \node[below] at (1.6,0) {\small runtime};
    \node[below] at (4.6,0) {\small hybrid};
    \node[below] at (7.6,0) {\small compile-time};

    \node[left] at (0,1.2) {\small ниско};
    \node[left] at (0,3.6) {\small средно};
    \node[left] at (0,6.0) {\small високо};

    \node[ormblock, minimum width=2.2cm] at (1.8,1.4) {Native API};
    \node[ormblock, minimum width=2.6cm] at (3.6,3.4) {Runtime query\\builder-и};
    \node[ormblock, minimum width=2.4cm] at (5.6,4.2) {Типизирани SQL\\DSL решения};
    \node[ormblock, minimum width=2.4cm] at (4.8,5.6) {Класически ORM\\+ кодогенерация};
    \node[ormaccent, minimum width=2.5cm] at (7.6,5.8) {Compile-time ORM};
\end{tikzpicture}
\caption{Позициониране на основните семейства решения според абстракция и момент на валидация}
\label{fig:existing_solutions_map}
\end{figure}

Фигура~\ref{fig:existing_solutions_map} показва, че compile-time ORM подходът не е просто ``още една ORM библиотека''. Той стои в горния десен квадрант, където високото ниво на абстракция се комбинира с ранна структурна проверка. Тази комбинация е рядка, защото изисква повече архитектурна дисциплина в самата библиотека.

\begin{table}[H]
\centering
\caption{Съпоставка на основните подходи за достъп до данни в C++}
\label{tab:existing_solutions_comparison}
\begin{tabular}{L{3.2cm}L{2.5cm}L{2.7cm}L{2.9cm}L{2.8cm}}
\toprule
\textbf{Подход} & \textbf{Ниво на абстракция} & \textbf{Валидация} & \textbf{Предимство} & \textbf{Основен недостатък} \\
\midrule
Native API & ниско & runtime & пълен контрол и близост до драйвера & силна обвързаност с конкретен backend \\
Runtime query builder & средно & основно runtime & по-четим API и по-малко boilerplate & заявката пак се материализира като низ \\
Типизиран SQL DSL & средно към високо & частично compile-time & по-добра статична структура за SQL света & обикновено остава ограничен до SQL \\
Класически ORM & високо & runtime или чрез кодогенерация & удобно object mapping поведение & сложна toolchain интеграция и непрозрачност \\
Compile-time ORM & високо & compile-time за структурата & ранно откриване на структурни грешки и формален connector contract & по-висока сложност на шаблонния слой \\
\bottomrule
\end{tabular}
\end{table}

\subsection{Критерии за сравнение}

За да бъде сравнението академично пълно, не е достатъчно да се каже, че една библиотека е ``по-удобна'', а друга е ``по-ниско ниво''. В настоящата работа се използват пет сравнителни критерия:

\begin{enumerate}[label=\arabic*.]
    \item \textbf{Ниво на логическа абстракция} --- доколко библиотеката позволява моделът на данните да се описва независимо от непосредствения синтаксис на backend-а.
    \item \textbf{Степен на статична валидация} --- кои аспекти на заявката и на модела могат да бъдат проверени преди изпълнение.
    \item \textbf{Преносимост между backend-и} --- възможно ли е логическият модел да бъде адаптиран към повече от един тип storage система.
    \item \textbf{Сложност на toolchain-а} --- изисква ли библиотеката външна кодогенерация, допълнителни preprocess стъпки или генератори на metadata.
    \item \textbf{Прозрачност на грешките и на изпълнението} --- доколко разработчикът може да проследи връзката между C++ кода, генерираната заявка и runtime поведението.
\end{enumerate}

Тези критерии са пряко свързани с целите на дипломната работа. Ако предложената библиотека трябва да бъде едновременно expressive и формално проверима, тогава тя трябва да превъзхожда алтернативите не само по удобство, а по баланса между тези пет измерения.

\subsection{Native клиентски библиотеки}

SQLite~\cite{sqlite}, \code{libpq} и \code{MySQL Connector/C} са представителни за подхода, при който приложението конструира и подава заявките директно към драйвера. Това е най-универсалният и често най-бързият вариант, но поставя тежестта на correctness изцяло върху приложния код. Разработчикът трябва ръчно да поддържа синхрон между имената на полета, поредността на параметрите и реалната схема на базата данни. При промяна на backend-а се променя не само диалектът на заявката, но и целият API за връзка, подготвяне, bind и четене на резултати.

Този подход е особено проблематичен, когато един и същ домейн модел трябва да бъде свързан последователно с повече от един тип база данни. В такъв случай приложението започва да натрупва условни пътища, backend-специфични SQL низове и дублирана логика за hydration на резултати. След определен размер на системата това намалява не само удобството, но и проверимостта на коректността.

Силната страна на native API подхода е неговата честност. Той не обещава повече абстракция, отколкото реално предоставя. Именно затова той често служи като контролна точка, спрямо която могат да бъдат оценявани по-високите abstraction слоеве.

\subsection{Runtime query builder-и}

Библиотеки като \code{SOCI}~\cite{soci} и част от по-лекия fluent SQL tooling търсят баланс между четимост и контрол. Те предлагат по-удобен C++ интерфейс за изграждане на заявки, но в много случаи продължават да разчитат на runtime описване и последващо рендиране към SQL. В практиката това означава, че известна част от структурата на заявката е по-добре капсулирана, но backend-независимата смяна на модел за съхранение остава ограничена.

За настоящата разработка тези библиотеки са важни като отправна точка. Те показват, че fluent API и типизирани елементи са полезни, но не са достатъчни сами по себе си. Същинският въпрос е дали компилаторът може да валидира цялата логическа структура на заявката и дали същата абстракция може да бъде пренесена извън SQL света.

\subsection{Типизирани SQL DSL решения}

\code{sqlpp11}~\cite{sqlpp11} е особено интересен случай, защото стои между традиционния runtime query builder и по-строго статичния дизайн. Подобни библиотеки използват шаблони и типове, за да направят част от SQL структурата видима за компилатора. Това е съществена стъпка напред спрямо чисто runtime builders, защото част от структурните грешки могат да бъдат хванати по-рано.

Същевременно обаче типизираният SQL DSL обикновено остава дълбоко свързан с релационния модел. Дори когато е много силен като C++ API, той най-често предполага свят, в който таблици, колони, joins и агрегати остават централната метафора. Това прави такива решения изключително ценни за SQL домейна, но по-трудно преносими към document, graph или key-value системи.

Именно тук се очертава една от основните разлики с настоящата разработка. Тя не цели просто compile-time структуриране на SQL. Тя цели compile-time логически слой, който може да бъде материализиран различно според backend семейството.

\subsection{Класически ORM рамки и кодогенерация}

\code{ODB}~\cite{odb} е характерен пример за ORM рамка, която предлага богато картографиране между класове и релационни таблици, но използва генерация на код и външен tooling. Това улеснява потребителя на ниво употреба, но измества сложността към build процеса и към автоматично генерираните артефакти. Получава се силна зависимост между модела на данните, допълнителната toolchain стъпка и конкретните релационни backend-и.

В класическите ORM рамки често се наблюдава и друго напрежение: колкото по-високо е удобството на ниво application код, толкова по-непрозрачен става пътят до реалната заявка и до performance характеристиките на системата. Това не е задължително фатален проблем, но в академичния контекст на настоящата работа е важно ограничение, защото затруднява проследимостта между абстракцията и материализацията.

\subsection{Граници на релационно центрираните абстракции}

За приложения, които трябва да работят с документни, графови или key-value системи, класическият ORM модел показва естествени ограничения. Част от концепциите --- като \code{JOIN}, външни ключове и строго таблична схема --- нямат директен еквивалент във всички модели за съхранение. Поради това обобщението на ORM отвъд релационния свят не може да бъде постигнато само чрез генериране на SQL.

Това наблюдение е важно и за native API, и за SQL DSL решенията. Дори когато те са много силни технически, обикновено предполагат, че SQL е централният език на света на данните. Настоящата разработка приема по-различна отправна точка: логическият модел трябва да бъде първичен, а релационната материализация --- само една от възможните интерпретации.

\begin{table}[H]
\centering
\caption{Ограничения на различните семейства решения спрямо multi-backend целта}
\label{tab:existing_solution_limitations}
\begin{tabularx}{\textwidth}{L{3.2cm}L{3.3cm}Y}
\toprule
\textbf{Семейство решения} & \textbf{Какво прави добре} & \textbf{Къде се появява ограничението} \\
\midrule
Native API & максимален контрол и близост до драйвера & липса на общ логически модел и голямо дублиране при повече от един backend \\
Runtime builders & по-добра четимост и капсулиране на заявката & валидацията остава предимно runtime и често е SQL-центрична \\
Типизирани SQL DSL & по-силна статична структура за SQL заявки & логиката обикновено остава ограничена в рамките на релационната метафора \\
Класически ORM & удобно object mapping поведение за релационни системи & зависимост от кодогенерация и трудна преносимост към нерелационни модели \\
\bottomrule
\end{tabularx}
\end{table}

\subsection{Позициониране на настоящата разработка}

Предложената библиотека се позиционира като \ORMName{} решение, което комбинира силна статична типизация, липса на външна кодогенерация и connector-базиран модел за пренасяне на логическата абстракция към различни backend-и. В сравнение с native API подхода тя намалява дублирането на логиката и риска от късно открити грешки. В сравнение с runtime query builder решенията премества структурната валидация към компилатора. В сравнение с класическите ORM рамки тя избягва външни генератори и поставя ясен акцент върху contract-а на конектора.

Ключовата разлика обаче е по-дълбока. Настоящата разработка не разглежда compile-time техниките като средство за ``по-умен SQL builder'', а като начин за формализиране на логическия модел на достъпа до данни. Точно това позволява следващите глави да преминат от теорията на типизираните модели към архитектурата, connector contract-а и реалните multi-backend сценарии.

\subsection{Изводи от сравнителния анализ}

Сравнението със съществуващите решения води до три основни извода. Първо, native API подходът остава незаменим като baseline за контрол и performance, но не предлага удовлетворително ниво на логическа абстракция. Второ, runtime query builder-ите и типизираните SQL DSL решения подобряват работата със SQL, но не решават напълно проблема за backend-независим логически слой. Трето, класическите ORM рамки предлагат висока абстракция, но често за сметка на toolchain сложност и ограничена пригодност извън релационния свят.

Точно това позициониране оправдава по-детайлното разглеждане в следващата глава на два отделни, но свързани въпроса: кои езикови и теоретични механизми позволяват compile-time ORM подход и как един и същ логически модел може да бъде интерпретиран в релационни и нерелационни системи.
